\section{Bernulli qonuni}

\begin{enumerate}

\item Gorizontal nayning ingichka qismidan oqayotgan suvning tezligi 2 m/s. Suvni nayning yug‘on qismidan oqish tezligini toping (m/s). Nayning ingichka va yug‘on qismlarining ko‘ndalang kesim yuzlari mos ravishda 200 $cm^2$ va 800 $cm^2$.
\begin{enumerate}[label=\Alph*)]
\item 0,1
\item 5
\item 0,5
\item 0,2
\end{enumerate}

\item Trubaning keng qismida suvning oqish tezligi 10 sm/s. Uning diametri keng qismiga qaraganda 4 marta kichik bo‘lgan qismida suvning oqish tezligi (m/s) qanday bo‘ladi?
\begin{enumerate}[label=\Alph*)]
\item 1,8
\item 1,6
\item 1,5
\item 1,7
\end{enumerate}

\item Chizmada tasvirlangan quvurlardan to‘la holda suyuqlik oqmoqda. Uning qaysi qismida suyuqlikning tezligi eng katta bo‘ladi?
\begin{enumerate}[label=\Alph*)]
\item 4
\item 3
\item 2
\item 1
\end{enumerate}

\item Rasmda gorizontal joylashgan trubadan suv to‘la holda oqayotgani ko‘rsatilgan. $\vartheta_1$ tezlik nimaga teng?
\begin{enumerate}[label=\Alph*)]
\item 4
\item 3
\item 2
\item 1
\end{enumerate}

\item Truba kesimining keng qismida suv 20 m/s tezlik bilan oqmoqda. Trubaning diametri 1,5 marta kichik bo‘lgan qismida oqim tezligi qanday bo‘ladi (m/s)?
\begin{enumerate}[label=\Alph*)]
\item 30
\item 45
\item 15
\item 12,6
\end{enumerate}

\item Ko‘ndalang kesim yuzi o‘zgaruvchi quvurdan suyuqlik oqmoqda (rasmga q.). Bu suyuqlikning qaysi kesimdagi bosimi eng kichik?
\begin{enumerate}[label=\Alph*)]
\item 3
\item 1
\item 2
\item oqayotgan suyuqlik qo‘shimcha bosim hosil qilmaydi
\end{enumerate}

\item Ko‘ndalang kesimi turli ($S_1; S_2; S_3$) bo‘lgan naylar orqali oqayotgan suv bosimlarining nisbati qaysi javobda to‘g‘ri keltirilgan?
\begin{enumerate}[label=\Alph*)]
\item $P_1 = P_2 = P_3$
\item $P_1 = P_2 < P_3$
\item $P_1 > P_2 > P_3$
\item $P_1 < P_2 < P_3$
\end{enumerate}

\item Suyuqlik ko‘ndalang kesimi o‘zgaruvchi nay bo‘ylab oqmoqda. Qaysi kesimda suyuqlik hosil qilayotgan bosim minimal bo‘ladi?
\begin{enumerate}[label=\Alph*)]
\item 1
\item 2
\item 3
\item 4
\end{enumerate}

\item Diametri 0,5 m bo‘lgan neft quvuri orqali 1 m/s tezlikda neft oqayotgan bo‘lsa, 1 soatda qancha miqdor neft olinadi (kg)? $\rho = 890 kg/m^3$.
\begin{enumerate}[label=\Alph*)]
\item $54,6 \cdot 10^5$
\item $4,5 \cdot 10^4$
\item $5,9 \cdot 10^5$
\item $5,1 \cdot 10^7$
\end{enumerate}

\item Trubadagi suvning tezligi 3,6 km/soat. Trubaning ko‘ndalang kesimidan har soatda 7,2 tonna suv oqib o‘tsa, trubaning ko‘ndalang kesimi yuzi ($sm^2$) qanchaga teng bo‘ladi?
\begin{enumerate}[label=\Alph*)]
\item 2
\item 20
\item 1
\item 10
\end{enumerate}

\item Katta shlankadan suv otilib chiqmoqda. Suv sarfi har minutda 60 l ga teng. Suvning 2 metr balandlikdagi ko‘ndalang kesim yuzasi ($sm^2$) nimaga teng. Shlang ichidagi yuza 1,5 $sm^2$ ga teng.
\begin{enumerate}[label=\Alph*)]
\item 3,32
\item 4,88
\item 4,37
\item 4,1
\end{enumerate}

\item Neft quduqdan diametri 60 mm bo‘lgan quvur orqali ko‘tariladi. Har soatda 9,12 t neft ko‘tarilayotgan bo‘lsa, neftning oqish tezligi topilsin (m/s). Neftning zichligi 800 $kg/m^3$.
\begin{enumerate}[label=\Alph*)]
\item 2,4
\item 1,2
\item 1,12
\item 6,2
\end{enumerate}

\item Ichida suvi bo‘lgan silindrik idishning tubiga diametri 2 sm bo‘lgan teshik ochildi va idishga har sekundda 360 g suv quyila boshlandi. Idishdagi suvning dastlabki balandligi (sm) qanday bo‘lganida uning sathi suv quyilib tursa o‘zgarmay turadi?
\begin{enumerate}[label=\Alph*)]
\item 7,2
\item 3,6
\item 6,7
\item 2,8
\end{enumerate}

\item Vodoprovod quvuri teshigidan tik yuqoriga otilayotgan suv oqimi 1,25 m ga ko‘tarilayotgan bo‘lsa, suv teshikdan necha m/s tezlik bilan otilib chiqmoqda?
\begin{enumerate}[label=\Alph*)]
\item 1,25
\item 2,5
\item 5,5
\item 5
\end{enumerate}

\item Daryo suvining tezligini o‘lchash maqsadida oqayotgan daryo suviga $90^\circ$ ga egilgan truba tushirilgan. Suv trubaga daryo suvining erkin sirtiga qaraganda 30 sm ga ko‘tarilgan bo‘lsa, daryo suvining tezligini (m/s) toping.
\begin{enumerate}[label=\Alph*)]
\item 2
\item $\sqrt{5}$
\item $\sqrt{3}$
\item $\sqrt{6}$
\end{enumerate}

\item Suv katerining tumshuq qismida vertikal bo‘ylab 1 m balandlikka ko‘tarilayotgan bo‘lsa, katerning tezligini taxminan qanday baholash mumkin (m/s)?
\begin{enumerate}[label=\Alph*)]
\item 8,2
\item 4,4
\item 1,6
\item 3,8
\end{enumerate}

\item Yuzasi 0,5 $sm^2$ bo‘lgan shpris porshenga 40 N kuch bilan bosim bersa, shpris teshigidan suv qanday (m/s) tezlik bilan oqib chiqadi?
\begin{enumerate}[label=\Alph*)]
\item 20
\item 10
\item 40
\item 15
\end{enumerate}

\item Agar suyuqlikning oqish tezligi 2 marta ortsa, uning perpendikulyar uriladigan to‘siqqa ko‘rsatadigan bosimi necha marta ortadi?
\begin{enumerate}[label=\Alph*)]
\item 2 marta
\item 4 marta
\item 8 marta
\item o‘zgarmaydi
\end{enumerate}

\item Agar bo‘yoq pultidan 25 m/s tezlik bilan suyuq bo‘yoq oqib chiqayotgan bo‘lsa, kompressor bo‘yoq pultida qanday bosim (Pa) hosil qiladi? Bo‘yoqning zichligi 800 $kg/m^3$.
\begin{enumerate}[label=\Alph*)]
\item $2 \cdot 10^5$
\item $1,5 \cdot 10^5$
\item $2,5 \cdot 10^5$
\item $10^5$
\end{enumerate}

\item Suv kraniga quvur ulangan. Kran 2 kPa bosim hosil qilayotgan bo‘lsa, quvurdan chiqayotgan suvning oqim tezligini toping (m/s).
\begin{enumerate}[label=\Alph*)]
\item 3
\item 1,2
\item 4,2
\item 2
\end{enumerate}

\item Purkagichdan suyuqlik 25 m/s tezlik bilan otilib chiqmoqda, suyuqlikning zichligi 1 $g/cm^3$. Purkagich idishining hajmida kompressor qanday bosim hosil qiladi (atm)?
\begin{enumerate}[label=\Alph*)]
\item 6,5
\item 3,1
\item 2,5
\item 3,9
\end{enumerate}

\item Ko‘ndalang kesim radiusi 4 sm bo‘lgan quvurdan zichligi 0,9 $g/cm^3$ bo‘lgan moy oqmoqda. Oqim tezligi 4 m/s ga teng. Yog‘ning harakat yo‘nalishidagi gidrodinamik bosimini toping (kPa).
\begin{enumerate}[label=\Alph*)]
\item 7,2
\item 4
\item 8,1
\item 1,2
\end{enumerate}

\item Diametri 12 sm bo‘lgan trubada suv 12 m/s tezlik bilan harakatlanmoqda, suv oqimining quvvatini toping (Wt).
\begin{enumerate}[label=\Alph*)]
\item 880
\item 8000
\item 3110
\item 9767
\end{enumerate}

\item Diametri 25 sm bo‘lgan trubada zichligi 800 $kg/m^3$ bo‘lgan neft qanday tezlik (m/s) bilan harakatlanganda, oqimining quvvati 2453,25 Wt bo‘ladi?
\begin{enumerate}[label=\Alph*)]
\item 3
\item 12
\item 6
\item 5
\end{enumerate}

\item Diametri 20 sm bo‘lgan teshikdan 4 m/s tezlik bilan chiqayotgan suv oqimining quvvatini toping (kVt).
\begin{enumerate}[label=\Alph*)]
\item 10
\item 1
\item 50
\item 1000
\end{enumerate}

\end{enumerate}

    

\begin{table}[h!]
\centering
\begin{tabular}{|c|c|c|c|c|c|c|c|c|c|}
\hline
№ & J. & № & J. & № & J. & № & J. & № & J. \\ \hline
1 & C & 7 & D & 11 & C & 16 & B & 22 & A \\ \hline
2 & B & 8 & B & 12 & C & 17 & C & 23 & D \\ \hline
3 & A & 9 & C & 13 & C & 18 & B & 24 & D \\ \hline
4 & B & 10 & B & 14 & D & 19 & C & 25 & B \\ \hline
5 & B & & & 15 & D & 20 & D & & \\ \hline
6 & C & & & & & 20 & B & & \\ \hline
\end{tabular}
\end{table}